% !PRINTED 392
\weberpage{392}
2. \gesperrt{If $p$ is an odd prime dividing $D$, then $a$ is a norm residue then and only then, when $a$ is a quadratic residue of $p$.}

This follows immediately from the condition for the norm residues:

\weq{4}{x^2 - D y^2 \equiv 4a \pmod{p}},

which in this case reduces to $x^2 \equiv 4a \pmod{p}$.

The prime 2 does not come into consideration, because $x^2 \equiv a \pmod{2}$ is always solvable; the moduli 4 and 8, however, do. If $D$ is odd, then $D \equiv 1 \pmod{4}$, and the condition for a norm residue $a$ is:

\[ \frac{x^2 - D y^2}{4} \equiv a \pmod{4} \text{ or } \pmod{8}, \]

and this congruence is solvable for every odd $a$ (by even $x, y$), hence:

3. \gesperrt{If 4 does not divide $D$, then every odd number $a$ is a norm residue of 4 and of 8.}

If $D$ is divisible by 4, then the condition for a norm residue $a$ is the possibility of the congruence:

\weq{5}{x^2 - \frac{D}{4} y^2 \equiv a \pmod{4} \text{ or } \pmod{8}.}

Since $a$ is odd, $x$ and $D y^2/4$ cannot both be odd.

If one goes through the individual cases accordingly, there results

\[
\begin{array}{l}
\dfrac{D}{4} \equiv 0 \pmod{8}, \quad a \equiv 1 \pmod{4}, \quad a \equiv 1 \pmod{8},\\[6pt]
\dfrac{D}{4} \equiv 1 \pmod{8}, \quad \text{no condition for } a,\\[6pt]
\dfrac{D}{4} \equiv 2 \pmod{8}, \quad a \equiv +1,\ -1 \pmod{8},\\[6pt]
\dfrac{D}{4} \equiv 3 \pmod{8}, \quad a \equiv 1 \pmod{4}, \quad a \equiv +1,\ 5 \pmod{8},\\[6pt]
\dfrac{D}{4} \equiv 4 \pmod{8}, \quad a \equiv 1 \pmod{4}, \quad a \equiv +1,\ 5 \pmod{8},\\[6pt]
\dfrac{D}{4} \equiv 5 \pmod{8}, \quad \text{no condition for } a,\\[6pt]
\dfrac{D}{4} \equiv 6 \pmod{8}, \quad a \equiv +1,\ 3 \pmod{8},\\[6pt]
\dfrac{D}{4} \equiv 7 \pmod{8}, \quad a \equiv 1 \pmod{4}, \quad a \equiv +1,\ 5 \pmod{8}.
\end{array}
\]
